% Section 19.10, from lectures/L19/appendix/c_exercises.md.

\section{Exercises}
\label{c19:sec:exercises}

\begin{aside}
Exercise~\ref{c19:ex:thirdnode}, \ref{c19:ex:feature} and \ref{c19:ex:afterabort} consolidate
\secref{c19:sec:verify}. Exercise~\ref{c19:ex:bridge} to \ref{c19:ex:endtoend} consolidate
\secref{c19:sec:bridge} to \ref{c19:sec:gaps} and appendix~\ref{app:protocol}, which is the authority
everywhere it and an exercise disagree.
\end{aside}

\exercise{A third node, and a new arbitration outcome}{Simulation}
\label{c19:ex:thirdnode}
\xpart{a} Extend \file{can_controller_tb.vhd} with a third node, C, sharing the same
\code{bus_line}. Give C the ID \code{0x002} and have it request transmission on the same cycle as node
A (\code{0x000}) and B (\code{0x001}). Confirm (by simulation, not just prediction) that A still
wins, and that \emph{both} B and C detect an arbitration loss.

\xpart{b} Node A's \code{0x000} is the lowest identifier a standard frame can carry, so no choice of
ID lets C beat it. Raise A's ID to \code{0x010}, and then choose an ID for C that makes C the winner
while B (\code{0x001}) still loses. Justify your choice with \secref{c10:sec:arbitration}'s
arbitration rule, then confirm the new winner in simulation. One consequence to notice before your
assertions surprise you: with A on \code{0x010}, B's \code{0x001} beats A as well, so in that round
\textbf{A loses too} and raises \code{error}; only C completes.

\exercise{A feature from a real controller, as a design sketch}{Reasoning}
\label{c19:ex:feature}
Pick \textbf{one} entry from \secref{c19:sec:gaps}'s list of what real CAN controllers do that this
design does not. Sketch, without necessarily implementing it in VHDL (in prose, plus an addition to
the state diagram or a new port list if that helps), what would have to change in
\code{can_controller} to support it. State which existing block or blocks (\code{bit_timer},
\code{crc15}, \code{tx_shift_reg}/\code{rx_shift_reg}, or \code{can_controller} itself) would have to
change, and which could stand exactly as they are.

\exercise{\code{spi_reg_bridge}}{VHDL}
\label{c19:ex:bridge}
Write your own \file{spi_reg_bridge.vhd} from \secref{c19:sec:bridge} and
appendix~\ref{app:protocol}. Verify it:

\begin{shell}
cd bridge
ghdl -a --std=93 spi_def.vhd spi_reg_bridge.vhd spi_reg_bridge_tb.vhd
ghdl -e --std=93 spi_reg_bridge_tb
ghdl -r --std=93 spi_reg_bridge_tb --assert-level=error
\end{shell}

\xpart{a} Get all four cases to pass. Cases 3 and 4 are the two that make this module more than a
byte counter \textemdash{} read their comments before you start, not afterwards.

\xpart{b} Now break it deliberately. Remove every use of \code{ss_active} and count bytes modulo five
instead. Cases 1 and 2 still pass. Say which check in case 3 fails, quote the value the error message
reports, and explain where those bytes came from.

\xpart{c} A bridge ignoring the reserved command bits also passes cases 1 and 2. Explain why the
protocol specification puts the reserved-bit rule in \emph{this} layer rather than in
\code{register_bank}, given that the bank never sees more than a 4-bit index.

\xpart{d} The read path latches \code{reg_rdata} \textbf{once}, at the end of the command byte. Say
what \code{register_bank} would have to do instead if the bridge reread it for every data byte, and
why that would be the wrong place to push the rule down into the bank. Put the module back.

\exercise{Composing \code{can_spi_node}}{VHDL}
\label{c19:ex:node}
Write the top level from \secref{c19:sec:node}: \code{spi_slave}, \code{spi_reg_bridge},
\code{register_bank} and \code{can_controller}, plus a \code{meta_prev} for the reset.

\xpart{a} Write it, and run \code{can_spi_node_tb} against it. Four testbenches handed out already
cover everything \emph{inside} the file; say which four, and then say the one kind of fault
\code{can_spi_node_tb} can fail that none of the four can. (The clue is the positional port binding.)

\xpart{b} \code{can_spi_node_tb} holds 200~ns of quiet around every \code{SS} edge, where
appendix~\ref{app:protocol} requires at least 60~ns. It is therefore deliberately kind to the timing
requirements, and no testbench handed out presses them. Copy it, shrink the quiet period step by step
towards 60~ns, and find out at what value your node stops answering.

Then say \textbf{why} exactly there, counted in clock cycles, and where in \code{spi_slave} the limit
sits. Compare the measured value with the 60~ns the specification promises: do you have margin, or is
your node right at the edge? This is the one place in the book where you \emph{measure} a margin
rather than read one, and it is worth doing carefully.

\xpart{c} Count the \code{meta_prev} instances in the finished design and say what each one
synchronises. Then say why \code{can_spi_node} hands \code{can_controller} the \emph{raw}
\code{reset_n} rather than the \code{reset_s2_n} it gives every other block, and what would go wrong
if it synchronised it twice.

\xpart{d} \code{tx_bus} and \code{bus_en} leave the design as two separate ports rather than one wire.
Explain what has to happen between them and a physical wire shared with a second board, and what would
fail if you assigned \code{tx_bus} straight to a pin and tied two boards together.

\exercise{What the board showed}{Reasoning}
\label{c19:ex:board}
The bring-up ran in eight steps, in two halves: three on two DE0-CVs and five with the hat as master.
This exercise is about what they established and what they did not. It needs no hardware of your own.

\xpart{a} The steps are ordered so that each rules out a class of faults before the next depends on
it. For each of steps 1, 4 and 5 (single board, loopback on the hat, register echo), say the one thing
it proves that the step before did not. Then say which of the eight steps \code{can_spi_node_tb} had
\textbf{already} established in simulation, and which only hardware can establish.

\xpart{b} Two wiring faults to tell apart. Say which step fails first if MISO and MOSI are swapped,
and which fails first if \code{VDDIO2} was never connected, leaving \code{PORTC} unsupplied. Explain
why they fail at different steps, and what that difference lets you conclude without touching a meter.
Also say what \code{MVIO.STATUS} would have answered in the second case, and why that answer is faster
than the meter.

\xpart{c} The test program handed out is written to succeed: it holds \code{SS} low in good time,
never sends a sixth byte and never aborts mid-transaction. Say which of the protocol specification's
rules \textemdash{} the 60~ns around \code{SS}, \code{MISO} during the command byte, or the abort rule
\textemdash{} \textbf{none} of the eight steps would have revealed as broken, and why that one is
therefore the most expensive to have wrong when a driver you did not write meets the node.

\xpart{d} A single node can never make \code{rx_valid} pulse, however its pins are wired, which is why
the session put a frame across \textbf{two} boards. Say why a node can never receive its own frame
(the reason is \chapref{c17:ch}'s \code{role}), and what the second board therefore proves that
neither a loopback nor \code{can_controller_tb} can.

\xpart{e} The two boards run on two independent 50 MHz oscillators. \Secref{c19:sub:gaps-tb} calls the
shared clock in \code{can_controller_tb} the testbench's \textbf{largest} gap for exactly that reason.
Be precise about what changes: \code{bit_timer.resync} already does real work in simulation, so say
what the two independent clocks add that no shared-clock simulation can show, and roughly how far
apart the two oscillators could drift before a frame stopped arriving intact
(Exercise~\ref{c12:ex:drift} did that calculation; reuse it). Compare that limit with a typical
crystal's tolerance and say whether the demonstration really puts it to the test.

\xpart{f} \code{rx_bus} comes from outside the FPGA, and \code{can_spi_node} passes it on
unsynchronised to \code{can_controller} deliberately, since \code{can_controller} contains its own
\code{meta_prev} (\chapref{c12:ch}). Explain what could go wrong if that instance were removed, why
the failure would be considerably harder to diagnose than a miswired line, and why no amount of
simulation in this book would ever reproduce it.

\exercise{Following a \code{send()} end to end}{Reasoning}
\label{c19:ex:endtoend}
\xpart{a} Follow a single frame request all the way down. Start from the five bytes on MOSI writing
\code{0x1} to \code{TX_SEND}, and follow them through \code{spi_slave} (bytes),
\code{spi_reg_bridge} (a register write) and \code{register_bank} (a one-cycle \code{tx_req}) to
\code{can_controller}'s ports. Name what every layer hands on to the next. Then do the same upwards,
for how the driver finds out that the frame is done.

\xpart{b} Four layers are three more than \code{can_controller_tb} needs. For every boundary, name the
testbench handed out that pins it down, and say what would be untested if that testbench did not
exist.

\xpart{c} Polling is all this design offers: there is no interrupt output anywhere. Say what would
have to be added to give the driver one, from which existing signal it would be derived, and in which
layer it would have to live. Then explain why it would have been hopeless for a driver to poll a
\code{tx_done} lasting one cycle, and which module made sure it is not.

\xpart{d} A driver writes \code{TX_ID}, then \code{TX_DLC}, then \code{TX_SEND}, as three separate
transactions. Between the second and the third the bus goes idle and another node transmits. Say
whether anything the driver has written is at risk, and which rule in the register map makes your
answer true.

\exercise{Life after an abort}{Simulation}
\label{c19:ex:afterabort}
\Secref{c19:sub:gaps-tb} lists ``no node ever transmits again after an abort'' as something the system
testbench does not cover. This exercise covers it, and shows what \code{crc15}'s \code{clear} input
was for.

\xpart{a} Extend \file{can_controller_tb.vhd} with a third case. Give both nodes time to return to
idle after case 2 has run, that is, after node B has lost arbitration and abandoned its frame
mid-identifier, and then have \textbf{B} send a frame of its own (ID \code{0x123}, DLC 2, any payload)
with A receiving. Confirm in simulation that A receives it: A's \code{rx_valid} pulses, A's
\code{error} stays low, and the ID comes back correctly.

\xpart{b} Now remove the clear: change \code{crc_clear} to a constant \code{'0'} and rerun. Cases 1
and 2 still pass. Case 3 does not. Say which node raises \code{error}, at which state, and explain why
it is \emph{not} that node that has the bug.

\xpart{c} Work out what B's \code{crc15} register holds at the moment it aborts in case 2, and why no
amount of correct behaviour afterwards ever returns it to zero. (The clue is \chapref{c13:ch}'s
property that generating and checking are the same operation: what does it assume about the register
at a frame's \emph{beginning}?)

\xpart{d} How many frames must node B send before it works again, with the clear still removed?
Explain your answer, and say what a driver polling \code{STATUS} and \code{ERROR_FLAGS} over SPI would
observe each time, using \chapref{c18:ch}'s rule for how the error bit is latched.

\xpart{e} \Chapref{c16:ch} puts the clear in \code{STATE_START} rather than in \code{STATE_CRC_WAIT} or
on the way out of an abort. Give one reason each why the other two placements are worse.
